\begin{abstract}
\noindent
How does road-network accessibility to resolutive health facilities vary across Tumbes, Amazonas and Cusco, Peru? We route \SampleSize{} sampled settlements to the nearest resolutive facility over an OpenStreetMap-derived road network and weight results by Census 2017 population (INEI, via CENEPRED), calibrated by district. The frame captures \FrameCoveragePct\% of census population -- far higher than its \PopulationMatchedPct\% settlement-count match rate alone suggests. Population-weighted mean car access time is \MeanAccessTotal{} minutes overall; \CovLeThirtyTotal\% of this population is within 30 minutes, while \CovNoTimeTotal\% has no estimable route -- a category kept distinct from ``\,$>$120 minutes.'' Access is markedly better in Tumbes (\MeanAccessTumbes{} min) than in Amazonas (\MeanAccessAmazonas{} min) or Cusco (\MeanAccessCusco{} min); the population-weighted Gini coefficient of access time is \GiniCalibrated. \TopOneDistrict{} (Cusco) is the largest computable gap, at \TopOneTime{} minutes. Findings are descriptive; no causal claims are made.
\end{abstract}
